\chapter{Installation and layout}
\label{ch:install}

\section{Requirements}

The trainer uses 64-bit floating point for every trainable tensor and every PDE derivative.
\file{swe_pinn/physics/quantum.py} enables JAX double precision when the package is imported, so JAX is imported even for a classical run.

Install the third-party packages from the repository root:

\begin{lstlisting}[style=sh]
pip install tensorflow optuna scipy matplotlib numpy pennylane jax jaxlib pandas
\end{lstlisting}

\begin{center}
\begin{tabular}{@{}lp{9.2cm}@{}}
\toprule
Package & Role \\
\midrule
TensorFlow / Keras &
Dense layers, Adam, \code{GradientTape} derivatives, saved weights \\
NumPy &
Collocation coordinates, exact-solution arrays, archives \\
Optuna &
Study, TPE sampler, SQLite storage, completed-trial cap \\
SciPy &
Optional L-BFGS-B; root finding for the dam-break and bump references \\
PennyLane &
Circuit definition and the \code{default.qubit} statevector device \\
JAX &
Batched just-in-time circuit and its vector--Jacobian products \\
Matplotlib &
Figures written by \file{analysis/analyze_study.py} \\
pandas &
Trial tables in the analysis script. The trainer does not import it \\
\bottomrule
\end{tabular}
\end{center}

PennyLane is imported only when a circuit is constructed.
A classical-only run still imports JAX, because the quantum module is imported with the package, and it does not construct a circuit.
A hybrid or full-quantum run needs both PennyLane and JAX.
The circuit is compiled with the JAX interface and differentiated with \code{diff_method="backprop"}.

Python~3.10 or newer matches the type hints in \file{user_input.py}.
The commands below use \code{python}; on some systems the interpreter is \code{python3}.
\file{run_cores.sh} calls \code{python3}.

\section{Check that the package imports}

From the repository root:

\begin{lstlisting}[style=sh]
python -c "from swe_pinn.physics.common import MODEL_FAMILIES; print(MODEL_FAMILIES)"
\end{lstlisting}

The expected line is

\begin{lstlisting}[style=sh]
('classical', 'hybrid_quantum', 'full_quantum')
\end{lstlisting}

A missing PennyLane installation does not fail this import.
It fails later, on the first hybrid or full-quantum model, with a message that names \code{pip install pennylane}.

\section{Repository layout}

Run every command from the repository root, the directory that contains \file{run_pinn.py} and \file{user_input.py}.

\begin{center}
\begin{tabularx}{\textwidth}{@{}lX@{}}
\toprule
Path & Role \\
\midrule
\file{user_input.py} &
The active input file. The solver reads \code{USER_PINN_INPUTS} and the physical dataclass selected by \code{problem_name}. \\
\file{user_input.py.*} &
Ready-made copies of that file for particular benchmarks and family sets. Copy one over \file{user_input.py} to use it. \\
\file{run_pinn.py} &
Process entry point. \\
\file{run_cores.sh} &
\(N\) search workers, then one final training. \\
\file{swe_pinn/cli.py} &
Argument parser for \code{--search-only}, \code{--worker-id}, and \code{--finalize}. \\
\file{swe_pinn/core/solver.py} &
Loss, Adam, L-BFGS-B, Optuna, and archives. \\
\file{swe_pinn/physics/common.py} &
The \code{PINN} network, collocation grids, and the problem dataclass. \\
\file{swe_pinn/physics/quantum.py} &
PennyLane circuit and the TensorFlow bridge. \\
\file{swe_pinn/problems/} &
Benchmark builders and the registry in \file{__init__.py}. \\
\file{swe_pinn/plots/jcp.py} &
Plot helpers. The solver itself writes NumPy archives; the study figures are made by the analysis script. \\
\file{analysis/analyze_study.py} &
Tradeoff report, Optuna figures, and the analytic comparison. \\
\file{analysis/<database>/} &
One folder per SQLite file: the database, \file{data/}, and \file{figures/}. \\
\file{logs/} &
Worker logs and the shared trial-table counter used by a parallel search. \\
\bottomrule
\end{tabularx}
\end{center}

A relative SQLite URL such as \code{sqlite:///optuna_tide.db} is stored at \file{analysis/optuna_tide/optuna_tide.db}.
The folder name is the database file name without \file{.db}.
Weights, \file{best_hyperparameters.txt}, and the NumPy archives go in \file{analysis/optuna_tide/data/}, whether or not Optuna actually ran.
Point \code{storage_url} at a new file before a smoke test so a short run does not replace the weights of a finished study.
Absolute SQLite URLs are left unchanged.

\section{First command}

After the packages import, the shortest confirmation that the entry point is wired up is the help text:

\begin{lstlisting}[style=sh]
python run_pinn.py --help
\end{lstlisting}

That command does not train a model.
The first training run is Example~\ref{ex:classical} in Chapter~\ref{ch:examples}.
Give it a new \code{storage_url} and a modest \code{final_adam_epochs} before you launch a multi-hour study.
